\metatitle{Gamma-positive polynomials}

\metadescription{Gamma-positivity of palindromic polynomials, the Foata--Strehl action, examples from permutation statistics, polytopes, and Gal's conjecture.}


\section[gammaPositivity]{Gamma positivity}

See the survey by \name{Christos Athanasiadis} \cite{Athanasiadis2018}.

A polynomial $p(x)$ of degree $n$ with non-negative coefficients
is \defin{palindromic} (or \defin{symmetric}) if $x^n p(x^{-1}) = p(x)$.
A palindromic polynomial $p(x)$ is said to be \defin{gamma-positive} if it can be expressed as
\[
 p(x) = \sum_{i=0}^{\lfloor n/2 \rfloor} \gamma_i x^i (1+x)^{n-2i}, \qquad \gamma_i \geq 0.
\]
The \defin{gamma-polynomial} of $p(x)$ is $\gamma_p(x) \coloneqq \sum_j \gamma_j x^j$.
If $p(x)$ is palindromic with center of symmetry $n/2$, the expansion above is unique.

Gamma-positivity immediately implies unimodality: writing $p(x) = \sum_k a_k x^k$ with
$a_k = a_{n-k}$ and expanding the binomial coefficients, each $\gamma_i$ contributes a
unimodal symmetric summand, so $p(x)$ is unimodal.

\begin{itemize}

\item $p(x)$ has all roots on the negative real axis if and only if $\gamma_p(x)$ does,
see \cite{Branden2004Poset,Gal2005}.
In particular, if $p(x)$ is real-rooted and palindromic, then it is gamma-positive
and $\gamma_p(x)$ is also real-rooted.

\item \name{Świętosław Gal} conjectured \cite{Gal2005} that the $h$-polynomial of any
flag simplicial sphere is gamma-positive.
This is still open in general but known for several families.

\item Suppose $p(x)$ is palindromic with center of symmetry $n/2$.
If $\gamma_p(x)$ is ultra-log-concave of order $\lfloor n/2 \rfloor$ with no internal zeros, then
$p(x)$ is also ultra-log-concave of order $n$ with no internal zeros, \cite{BrandenFerroniJochemko2024x}.
The converse does not hold.

\item Suppose $p(x)$ is palindromic.
If $\gamma_p(x)$ is log-concave with no internal zeros, then
$p(x)$ is also log-concave with no internal zeros, \cite{FerroniPanovaVenturello2025x}.
The converse does not hold.

\end{itemize}

\todo{Add https://arxiv.org/pdf/2302.00754.pdf }
\todo{https://arxiv.org/pdf/2202.08984.pdf}
\todo{https://arxiv.org/pdf/2306.15785.pdf}


\subsection[gammaPositiveExamples]{Examples of gamma-positivity}


\begin{theorem}[Foata--Strehl action on Eulerian polynomials]
The \hyperref[interlacingExamples]{Eulerian polynomials} $A_n(x)$ are palindromic and gamma-positive.
The gamma-coefficients $\gamma_{n,k}$ count permutations $\sigma \in \symS_n$
with $k$ \defin{peaks} (positions $i$ with $\sigma_{i-1} < \sigma_i > \sigma_{i+1}$).
Thus
\[
 A_n(x) = \sum_{k=0}^{\lfloor (n-1)/2 \rfloor} \gamma_{n,k}\, x^{k+1}(1+x)^{n-1-2k}.
\]
The proof uses the \defin{Foata--Strehl} (or \emph{valley-hopping}) action on $\symS_n$:
for each $i \in \{2,\dotsc,n-1\}$, define an involution $\varphi_i$ that swaps the
relative position of $i$ with its neighbors when $i$ is a local minimum.
The $2^{n-2}$ orbits of the abelian group generated by $\{\varphi_i\}$ are exactly the
gamma-classes; orbits of size $2^k$ contribute $\gamma_{n,k}$.
See \cite{Athanasiadis2018} and \cite{Branden2008} for generalization to sign-graded posets.
We also get gamma-positivity for the
\hyperref[interlacingExamples]{peak polynomials} by the same method.
\end{theorem}


\begin{theorem}[Narayana polynomials, \cite{Branden2006}]
The \hyperref[realRootedCatalan]{Narayana polynomials}
\[
 N_n(x) = \frac{1}{n}\sum_{k=1}^n \binom{n}{k}\binom{n}{k-1} x^k
\]
are palindromic and gamma-positive. The gamma-coefficients are the Narayana numbers of the second kind.
Since $N_n(x)$ is also real-rooted (see \cite{Branden2006}), the gamma-positivity follows
from the real-rooted criterion above.
\end{theorem}


\begin{theorem}[Descent polynomials of standard Young tableaux]
Let $W_\lambda(t) = \sum_{T \in \SYT(\lambda)} t^{\des(T)}$.
For any shape $\lambda$, $W_\lambda(t)$ is palindromic and gamma-positive.
This follows because $W_\lambda(t)$ is real-rooted
(see \cite{Brenti1989} and the discussion in \hyperref[realRootedTableaux]{tableaux examples}),
combined with the criterion that palindromic real-rooted polynomials are gamma-positive.
\end{theorem}


\begin{theorem}[Zig-zag poset, \cite{PetersenZhuang2024x}]
Let $Z_m(t)$ be the $P$-Eulerian polynomial associated with the zig-zag poset on $m$ vertices.
Then $Z_m(t)$ is gamma-positive:
\[
 Z_m(t) = t \cdot \sum_{0 \leq 2j \leq m-2} \gamma_{m,j}\,t^j (1+t)^{m-2-2j}, \qquad \gamma_{m,j} \geq 0.
\]
We have that
\[
 Z_m(t) = t \cdot \sum_{\pi \in U_m} t^{\ret_1(\pi)}
\]
where the sum is over \hyperref[namedPermutationSets]{up-down permutations},
and $\ret_1(\pi)$ counts instances where $i+1$ appears to the left of $i$ but not adjacent to it.
It is an open problem to give a combinatorial explanation for the symmetry and unimodality
of the $\gamma_{m,j}$.
\end{theorem}


\begin{theorem}[Brändén, \cite{Branden2004Poset}]
For any graded poset $P$ with a $\hat{0}$ and $\hat{1}$,
if the \defin{$cd$-index} of $P$ has non-negative coefficients,
then the $h$-polynomial of the order complex of the proper part $\overline{P}$
is gamma-positive.
In particular, this applies to face posets of polytopes, Boolean algebras,
and the partition lattice.
\end{theorem}


\subsection[gammaPositivityMethods]{Methods for proving gamma-positivity}

There are several standard approaches:

\begin{enumerate}

\item \emph{Real-rootedness.}
If $p(x)$ is palindromic and real-rooted with non-negative coefficients,
then it is automatically gamma-positive.
This gives gamma-positivity for Eulerian, Narayana, and SYT descent polynomials.

\item \emph{Group action.}
Exhibit an involution (or group action) whose orbits correspond to the gamma-classes.
The Foata--Strehl action is the canonical example.

\item \emph{Symmetric decomposition.}
Write $p(x) = a(x)(1+x) + x^c b(x)$ where $a,b$ are palindromic;
if $a$ and $b$ are gamma-positive, so is $p$.
See \cite{BrandenSolus2019} for this approach via the symmetric decomposition of the $h^*$-polynomial.

\item \emph{Algebraic / representation-theoretic.}
Show that $p(x)$ is the Hilbert series of a representation where the palindromy
and gamma-positivity are forced by the structure, e.g.\ via
the coinvariant algebra or $\mathfrak{sl}_2$-module structure.

\end{enumerate}


\section[realrootedOpenProblems]{Open problems}

In \cite{CoulterDoMoskovsky2023x}, the authors consider polynomials related to Hurwitz numbers.
These are conjectured to be real-rooted and interlace in a certain fashion \cite[Conj. 3.14]{CoulterDoMoskovsky2023x}.
The authors prove this statement for one-point Hurwitz numbers.

Gal's conjecture \cite{Gal2005} that $h$-polynomials of flag simplicial spheres are gamma-positive
remains open in general.





\section[symmetricDecomposition]{Symmetric decomposition}

See \cite{BrandenSolus2019} for the symmetric decomposition
and relationship with real-rootedness.
